\chapter{幾何力学入門：接続・曲率・ホロノミー}\label{ch:basics-geom}

AMBAC、落下する猫、宇宙飛行士の姿勢変更、コッドマンのパラドックスは、すべて同じ数学で説明できる。
本章では、その数学（幾何力学）の考え方を、式の厳密さよりも直観を優先して紹介する。詳しくは \cite{R5,R4,R6} を参照。

\section{形状と位置・姿勢を分ける}

自由に浮かぶ機体の状態は、次の 2 つに分けて考えると見通しがよい。
\begin{itemize}
\item \textbf{形状}（shape）：関節角 $q$。機体の「形」で、胴体から見た手足の配置。
\item \textbf{位置・姿勢}（group）：重心の位置と胴体の向き $\theta$（3 次元なら $R$）。形を保ったまま機体全体を動かしても変わらない量。
\end{itemize}
数学的には、全体の配位空間は「形状空間」の各点に「位置・姿勢の空間」が付いた構造（主ファイバー束）を持つ。
物理法則（運動エネルギー）が位置・姿勢の取り方によらない（対称性がある）ことが、以下の議論の前提になる。

\section{接続：形を変えたとき向きがどう変わるか}

角運動量がゼロのとき、形を少し変える（$dq$）と、向きは決まった量だけ変わる：
\begin{equation}
d\theta = -\Gamma(q)\cdot dq .
\end{equation}
この「形の変化から向きの変化を決める規則」$\Gamma$ を\textbf{接続}（connection）という。力学から決まる接続なので、機械的接続と呼ぶ。

\begin{keypoint}
接続は、形状空間の各点に「どちらに形を変えると、向きがどれだけ変わるか」を与える。時間や速さは含まれない。
したがって、形の変え方（経路）だけで向きの変化が決まる。
\end{keypoint}

\paragraph{例：手足 1 本の平面モデル}
形状空間は 1 次元（$q$ だけ）で、$\Gamma(q) = B(q)/A(q)$。付録~\ref{app:twobody} の閉じた形を使えば、腕の角度ごとの $\Gamma$ を手計算でも求められる。

\section{曲率とホロノミー}

\subsection{閉じた経路で何が起きるか}

形状を変えて元に戻す（形状空間の閉曲線 $C$ を 1 周する）と、向きの変化は
\begin{equation}
\Delta\theta = -\oint_C\Gamma\cdot dq .
\end{equation}
形状空間が 1 次元なら、閉曲線は「行って戻る」だけなので $\Delta\theta=0$ である。2 次元以上なら、閉曲線は面積を囲めるので、一般に $\Delta\theta\neq0$ となる。
この「1 周したときに残るずれ」を\textbf{ホロノミー}（holonomy）、あるいは\textbf{幾何学的位相}（geometric phase）という。

\subsection{曲率}

ストークスの定理により、閉曲線に沿った積分は、囲まれた領域での「曲率」の面積分に置き換えられる：
\begin{equation}
\Delta\theta = -\iint_S F\,dq_1\,dq_2,\qquad F = \frac{\partial\Gamma_2}{\partial q_1}-\frac{\partial\Gamma_1}{\partial q_2}.
\end{equation}
曲率 $F$ は「形状空間の各点で、小さなループを回したときにどれだけ向きが進むか」の密度である。
曲率マップ（第~\ref{ch:planar}章、第~\ref{ch:rw}章）は、この密度を色で表したもので、AMBAC の「効きやすい形の領域」を示す地図になる。

\begin{example}[小さな正方形ループ]
形状空間で一辺 $\varepsilon$ の小さな正方形を反時計回りに回すと、
4 辺の積分は、向かい合う辺どうしで $\Gamma$ の差を拾うので
\[
\Delta\theta \approx -\left[\Gamma_1(q_2)-\Gamma_1(q_2+\varepsilon)\right]\varepsilon - \left[\Gamma_2(q_1+\varepsilon)-\Gamma_2(q_1)\right]\varepsilon
\approx -F\,\varepsilon^2 .
\]
一辺を 2 倍にすると回転は 4 倍になる（2 次の効果）。Sim02 の曲率マップで $|F|\approx0.04$ の領域に一辺 0.5 rad（約 29°）の正方形を置くと、
1 周あたり $0.04\times0.25 = 0.01$ rad（約 0.6°）の回転が得られる。
\end{example}

\subsection{可換と非可換}

平面では向き $\theta$ が 1 つの数なので、ループを続けて回すと回転は足し算になり、順序によらない（可換）。
3 次元では向きが回転行列なので、ループごとの回転を合成すると順序によって結果が変わる（非可換）。
このとき曲率には、$\Gamma$ の微分に加えて、$\Gamma$ どうしの交換子の項が加わる。
非可換性は、第~\ref{ch:3d}章の「ヨー → ピッチ」と「ピッチ → ヨー」の違いとして現れた。

\section{非ホロノミック系とリー括弧}

\subsection{ホロノミックと非ホロノミック}

速度に関する拘束（ここでは $L=0$、すなわち $\dot\theta = -\Gamma\cdot\dot q$）が、位置に関する拘束（$\theta = f(q)$）に積分できるとき、ホロノミックという。
積分できないとき、非ホロノミックという。手足 1 本の AMBAC はホロノミック（$\theta$ は $q$ の関数）、手足 2 本以上で曲率がゼロでない AMBAC は非ホロノミックである。

非ホロノミック系の身近な例は、車の縦列駐車である。車は真横には動けない（速度の拘束）が、前後進とハンドル操作を組み合わせると、横にずれることができる。
AMBAC で「形を元に戻しながら向きだけを変える」のは、これと同じ構造である。

\subsection{リー括弧}

2 つの動き（ベクトル場）$X$, $Y$ を、$X$ を $\varepsilon$ → $Y$ を $\varepsilon$ → $X$ を $-\varepsilon$ → $Y$ を $-\varepsilon$ の順に行うと、
元の点には戻らず、$\varepsilon^2[X,Y]$ だけずれる。この $[X,Y]$ を\textbf{リー括弧}という。
リー括弧の方向にも動けることが、非ホロノミック系の可制御性の鍵である（Chow--Rashevskii の定理）。
AMBAC の曲率は、このリー括弧の向き成分にあたる。正弦波の入力で系を操る方法 \cite{R10} は、この考え方に基づく。

\subsection{Brockett の定理}

非ホロノミック系は、可制御であっても、滑らかで時間に依存しないフィードバック則では原点に漸近安定化できない \cite{R9}。
AMBAC で姿勢を目標に収束させるには、周期的な入力（ループ運動）、時間に依存する制御、あるいは切り替えが必要になる。
第~\ref{ch:ai}章で、構造のない学習方策が失敗し、振動子（周期的な構造）を持つ方策が成功したことは、この定理と整合している。

\section{同じ数学が現れる他の例}

\begin{longtable}{p{0.24\linewidth}p{0.30\linewidth}p{0.38\linewidth}}
\toprule
現象 & 形状空間 & 1 周で残るずれ \\ \midrule \endhead
AMBAC（本報告書） & 関節角 & 胴体の向き \\
落下する猫 \cite{R4,R7} & 上半身と下半身の曲げ・ひねり & 体の向き（足から着地できる） \\
宇宙飛行士の姿勢変更 \cite{R8} & 腕・脚の関節角 & 体の向き \\
コッドマンのパラドックス（第~\ref{ch:shoulder}章） & 腕の向き（単位球面） & 上腕のひねり（＝立体角） \\
低レイノルズ数の泳ぎ \cite{R6} & 微生物の体の形 & 位置（前進） \\
車の縦列駐車 & 前後進とハンドル角 & 横方向の位置 \\
\bottomrule
\end{longtable}

\section*{考察課題}
\begin{exercise}
形状空間で、曲率 $F$ が正の領域と負の領域にまたがる大きなループを回すと、1 周あたりの回転は小さくなる。その理由を説明せよ。
\end{exercise}
\begin{exercise}
腕と脚を同じ位相で同時に往復させた場合（形状空間で直線上を往復する場合）、正味の回転がゼロになることを示せ。
\end{exercise}
